\documentclass[aspectratio=169,12pt]{beamer}
\usepackage[spanish,es-nodecimaldot]{babel}
\usepackage[T1]{fontenc}
\usepackage[utf8]{inputenc}
\usepackage{tgheros}
\renewcommand{\familydefault}{\sfdefault}
\usepackage{csquotes}
\usepackage[backend=bibtex,style=ieee,sorting=none,maxbibnames=2,minbibnames=1,doi=true,url=false,isbn=false]{biblatex}
\addbibresource{referencias.bib}
% Las notas sobre las versiones consultadas permanecen en el archivo .bib.
\AtEveryBibitem{\clearfield{note}}
\urlstyle{same}
\definecolor{Paper}{HTML}{F7F5EF}
\definecolor{Night}{HTML}{112C36}
\definecolor{Ink}{HTML}{19343B}
\definecolor{Accent}{HTML}{00766E}
\definecolor{Amber}{HTML}{E7AD62}
\definecolor{Muted}{HTML}{5A6C70}
\definecolor{Pale}{HTML}{B7C9CA}
\setbeamercolor{normal text}{fg=Ink,bg=Paper}
\setbeamercolor{background canvas}{bg=Paper}
\setbeamercolor{structure}{fg=Accent}
\setbeamercolor{frametitle}{fg=Ink,bg=Paper}
\setbeamercolor{footline}{fg=Muted,bg=Paper}
\setbeamerfont{frametitle}{size=\Large,series=\bfseries}
\setbeamertemplate{navigation symbols}{}
\setbeamersize{text margin left=11mm,text margin right=11mm}
\setbeamertemplate{bibliography item}{\insertbiblabel}
\setbeamerfont{bibliography item}{size=\scriptsize}
\setbeamercolor{bibliography entry author}{fg=Ink}
\setbeamercolor{bibliography entry title}{fg=Ink}
\setbeamercolor{bibliography entry location}{fg=Ink}
\setbeamercolor{bibliography entry note}{fg=Ink}
\newcommand{\etiqueta}[1]{{\fontsize{7.5}{9}\selectfont\bfseries\color{Accent}#1\par}}
\newcommand{\texto}[1]{{\fontsize{11}{14}\selectfont #1\par}}
\newcommand{\fuentesactuales}{}
\newcommand{\fuentesdiapositiva}[1]{\gdef\fuentesactuales{#1}}
\setbeamertemplate{frametitle}{\vspace*{4mm}\etiqueta{\secname}\vspace{1.8mm}{\usebeamerfont{frametitle}\insertframetitle\par}\vspace{3mm}}
\setbeamertemplate{footline}{%
  \begin{beamercolorbox}[wd=\paperwidth,ht=9mm,dp=2mm,
    leftskip=11mm,rightskip=11mm]{footline}
    {\fontsize{6.8}{8}\selectfont
      \ifx\fuentesactuales\empty
        \strut
      \else
        \textbf{Fuentes}\hspace{2mm}\fuentesactuales
      \fi\par}
    \vspace{1.5mm}
    {\fontsize{6.5}{8}\selectfont
      ISR401-BCCH\hfill EXP-1\hspace{5mm}
      {\bfseries\insertframenumber\,/\,\inserttotalframenumber}\par}
  \end{beamercolorbox}
}
\title{IA generativa como apoyo a la elicitación}
\author{Equipo BCCH}
\date{}
\begin{document}
{
\setbeamercolor{background canvas}{bg=Night}
\setbeamercolor{normal text}{fg=white,bg=Night}
\begin{frame}[plain,t]
\vspace*{8mm}
{\fontsize{7.5}{9}\selectfont\bfseries\color{Pale}INGENIERÍA DE REQUERIMIENTOS\par}
\vspace{7mm}
\begin{columns}[T,onlytextwidth]
\begin{column}{0.76\textwidth}
{\fontsize{27}{30}\selectfont\bfseries\color{white}IA generativa\\como apoyo\\a la elicitación\par}
\vspace{4mm}
{\fontsize{12}{15}\selectfont\color{Pale}Capacidades, límites y riesgos\par}
\end{column}
\begin{column}{0.19\textwidth}
\vspace{2mm}
{\fontsize{7.5}{9}\selectfont\color{Pale}TEMA\par}
\vspace{1mm}
{\fontsize{48}{50}\selectfont\bfseries\color{Amber}08\par}
\end{column}
\end{columns}
\vfill
{\fontsize{9}{11}\selectfont\color{white}Equipo BCCH\par}
\vspace{1.5mm}
{\fontsize{7.5}{9}\selectfont\color{Pale}GA-SUM-02 · EXP-1\par}
\vspace{7mm}
\end{frame}
}
\section{FUNDAMENTOS}
\begin{frame}[t]{Elicitación de requisitos}
\fuentesdiapositiva{\cite{ireb2026foundation} IREB, 2026\qquad
  \cite{washizaki2024swebok} SWEBOK, 2024}
\begin{columns}[T,onlytextwidth]
\begin{column}{0.47\textwidth}
{\fontsize{23}{27}\selectfont\bfseries Comprender\\necesidades\par}
\vspace{4mm}
\texto{Hacer explícitas las expectativas de los interesados.}
\end{column}
\begin{column}{0.45\textwidth}
\etiqueta{FUENTES}
\texto{Personas, documentos y sistemas existentes.}
\vspace{2.5mm}
\etiqueta{TÉCNICAS}
\texto{Entrevistas, observación, talleres y prototipos, según el contexto.}
\vspace{2.5mm}
\etiqueta{CONOCIMIENTO}
\texto{Las necesidades también pueden ser implícitas o difíciles de expresar.}
\end{column}
\end{columns}
\vfill
{\fontsize{10}{13}\selectfont\bfseries\color{Accent}Los requisitos candidatos necesitan análisis y validación.\par}
\vspace{2mm}
\end{frame}
\begin{frame}[t]{IA generativa y modelos de lenguaje}
\fuentesdiapositiva{\cite{cheng2026genai} Cheng et al., 2026\qquad
  \cite{vogelsang2025guideline} Vogelsang y Fischbach, 2025}
\vspace{2mm}
\begin{columns}[T,onlytextwidth]
\begin{column}{0.46\textwidth}
\etiqueta{IA GENERATIVA}
\vspace{2mm}
{\fontsize{20}{24}\selectfont\bfseries Genera contenido\\con patrones\\aprendidos.\par}
\end{column}
\begin{column}{0.46\textwidth}
\etiqueta{LLM GENERATIVO}
\vspace{2mm}
\texto{Un modelo de lenguaje de gran tamaño puede generar texto a partir de instrucciones y contexto.}
\vspace{2mm}
\etiqueta{PROMPT}
\texto{La entrada expresa la tarea y puede aportar datos y un formato de respuesta.}
\end{column}
\end{columns}
\vfill
{\fontsize{10}{13}\selectfont\bfseries\color{Accent}Una respuesta fluida puede contener errores.\par}
\vspace{2mm}
\end{frame}
\section{APLICACIONES}
\begin{frame}[t]{Apoyo durante la elicitación}
\fuentesdiapositiva{\cite{cheng2026genai} Cheng et al., 2026\qquad
  \cite{arora2024advancing} Arora et al., 2024}
\vspace{1mm}
\begin{columns}[T,onlytextwidth]
\begin{column}{0.08\textwidth}{\fontsize{20}{24}\selectfont\color{Accent}01\par}\end{column}
\begin{column}{0.86\textwidth}{\large\bfseries Preparar preguntas\par}\texto{Explorar necesidades y aclarar aspectos del dominio.}\end{column}
\end{columns}
\vspace{5mm}
\begin{columns}[T,onlytextwidth]
\begin{column}{0.08\textwidth}{\fontsize{20}{24}\selectfont\color{Accent}02\par}\end{column}
\begin{column}{0.86\textwidth}{\large\bfseries Examinar información\par}\texto{Organizar documentos y señalar posibles vacíos o ambigüedades.}\end{column}
\end{columns}
\vspace{5mm}
\begin{columns}[T,onlytextwidth]
\begin{column}{0.08\textwidth}{\fontsize{20}{24}\selectfont\color{Accent}03\par}\end{column}
\begin{column}{0.86\textwidth}{\large\bfseries Proponer requisitos candidatos\par}\texto{Formular propuestas para contrastarlas con sus fuentes y validarlas con los interesados.}\end{column}
\end{columns}
\vfill
\vspace{2mm}
\end{frame}
\section{LÍMITES Y RIESGOS}
\begin{frame}[t]{Riesgos y revisión humana}
\fuentesdiapositiva{Riesgos: \cite{cheng2026genai,arora2024advancing}\qquad
  Criterios de revisión: \cite{ireb2026foundation,vogelsang2025guideline}}
\begin{columns}[T,onlytextwidth]
\begin{column}{0.46\textwidth}
\etiqueta{RIESGO}
\end{column}
\begin{column}{0.44\textwidth}
\etiqueta{VERIFICACIÓN}
\end{column}
\end{columns}
\vspace{3mm}
\begin{columns}[T,onlytextwidth]
\begin{column}{0.46\textwidth}
{\bfseries Información inventada}\par
\texto{Reglas sin respaldo.}
\end{column}
\begin{column}{0.44\textwidth}
\texto{Vincular cada requisito con su fuente.}
\end{column}
\end{columns}
\vspace{4mm}
\begin{columns}[T,onlytextwidth]
\begin{column}{0.46\textwidth}
{\bfseries Sesgos y contexto}\par
\texto{Perspectivas omitidas.}
\end{column}
\begin{column}{0.44\textwidth}
\texto{Contrastar propuestas con los interesados.}
\end{column}
\end{columns}
\vspace{4mm}
\begin{columns}[T,onlytextwidth]
\begin{column}{0.46\textwidth}
{\bfseries Privacidad}\par
\texto{Datos sin protección.}
\end{column}
\begin{column}{0.44\textwidth}
\texto{Resguardar los datos de los interesados.}
\end{column}
\end{columns}
\vfill
\vspace{1mm}
\end{frame}
\section{CASO DE PRÁCTICA}
\begin{frame}[t]{PrintFlow: pedidos de impresión}
\fuentesdiapositiva{}
\begin{columns}[T,onlytextwidth]
\begin{column}{0.47\textwidth}
{\fontsize{21}{25}\selectfont\bfseries Pedidos\\personalizados\par}
\vspace{3mm}
\texto{El caso plantea archivos y aprobaciones dispersos que dificultan
identificar el diseño autorizado.}
\vspace{3mm}
\etiqueta{OBJETIVO PROPUESTO}
\texto{Centralizar pedidos y aprobaciones para coordinar la producción.}
\end{column}
\begin{column}{0.45\textwidth}
\etiqueta{ACTORES DEL TALLER}
\vspace{2mm}
\texto{\textbf{Cliente}\quad Revisa el diseño y responde.}
\vspace{2mm}
\texto{\textbf{Vendedor}\quad Registra el pedido y sus artículos.}
\vspace{2mm}
\texto{\textbf{Operador}\quad Carga versiones y actualiza el avance.}
\vspace{2mm}
\texto{\textbf{Administrador}\quad Consulta pendientes y atrasos.}
\end{column}
\end{columns}
\end{frame}

\begin{frame}[t]{Alcance de la demostración}
\fuentesdiapositiva{}
{\fontsize{17}{21}\selectfont\bfseries Aprobación y paso a producción\par}
\vspace{3mm}
\begin{columns}[T,onlytextwidth]
\begin{column}{0.45\textwidth}
\etiqueta{INFORMACIÓN DISPONIBLE}
\vspace{2mm}
\texto{Versiones de diseño conservadas.}
\vspace{2.5mm}
\texto{Aprobación o solicitud de cambios.}
\vspace{2.5mm}
\texto{Control del paso a producción.}
\end{column}
\begin{column}{0.47\textwidth}
\etiqueta{PREGUNTAS POR ACLARAR}
\vspace{2mm}
\texto{¿Qué ocurre si se carga otra versión después de aprobar?}
\vspace{3mm}
\texto{¿La aprobación corresponde a un artículo o a todo el pedido?}
\end{column}
\end{columns}
\vfill
{\fontsize{9}{11}\selectfont\color{Accent}\bfseries
La IA propondrá preguntas y requisitos candidatos para revisión humana.\par}
\vspace{1mm}
\end{frame}

\section{DEMOSTRACIÓN}
\begin{frame}[t]{Prueba con IA en PrintFlow}
\fuentesdiapositiva{}
\begin{columns}[T,onlytextwidth]
\begin{column}{0.44\textwidth}
\etiqueta{ENTRADA ENVIADA}
\vspace{1mm}
{\fontsize{10}{12.5}\selectfont Instrucción y contexto completo del caso, en un solo mensaje.\par}
\vspace{2mm}
\etiqueta{EJECUCIÓN}
\vspace{1mm}
{\fontsize{10}{12.5}\selectfont Gemini, versión gratuita. Una ejecución, 21 de noviembre de 2026.\par}
\end{column}
\begin{column}{0.48\textwidth}
\etiqueta{FRAGMENTO DE LA RESPUESTA ORIGINAL (RC-04)}
\vspace{1mm}
{\fontsize{9.5}{12}\selectfont «El sistema impedirá que un artículo pase a producción o que el Operador avance su estado de producción si no existe una aprobación vigente correspondiente a la versión de diseño actual del artículo.»\par}
\vspace{1mm}
{\fontsize{7.5}{9}\selectfont\color{Muted}Respuesta completa: EXP-1/evidencias/respuesta-ia.txt\par}
\end{column}
\end{columns}
\vspace{2mm}
\etiqueta{QUÉ PRODUJO}
\vspace{1mm}
{\fontsize{10}{12.5}\selectfont 11 preguntas distribuidas entre los actores y 9 requisitos candidatos (RC-01 a RC-09), más 6 suposiciones y 6 decisiones por confirmar.\par}
\vfill
{\fontsize{9}{11}\selectfont\color{Accent}\bfseries
Propuestas de una sola ejecución, pendientes de revisión humana.\par}
\vspace{1mm}
\end{frame}

% PARTE 2 — Revision humana de la respuesta de IA
\section{REVISIÓN HUMANA}
\begin{frame}[t,shrink=14]{Revisión humana de la respuesta de IA}
\fuentesdiapositiva{PrintFlow: contexto F01--F05 y evidencias/respuesta-ia.txt}
\vspace{1mm}
\begin{columns}[T,onlytextwidth]
\begin{column}{0.47\textwidth}
\etiqueta{HALLAZGOS COMPROBABLES}
\vspace{1.5mm}
{\fontsize{10}{13}\selectfont
\textbf{9 requisitos candidatos} (RC-01 a RC-09).\\[2mm]
\textbf{8/9} citan F01--F05; RC-09 se apoya en la descripción del Administrador, fuera de esos hechos.\\[2mm]
\textbf{9/9} requieren confirmar al menos una regla.\\[2mm]
\textbf{11 preguntas}: 6 claras sin cambios y 5 con mejoras de redacción sugeridas.
\par}
\end{column}
\begin{column}{0.47\textwidth}
\etiqueta{EJEMPLOS REALES DE LA RESPUESTA}
\vspace{1.5mm}
{\fontsize{9.3}{12}\selectfont
\textbf{RC-02:} «dejará sin efecto cualquier aprobación previa»; F02/F03/F05 identifican versiones y riesgo, \emph{no confirman} revocación automática.\\[2mm]
\textbf{RC-04:} «impedirá que un artículo pase a producción»; relacionado con F04/F05, pero falta confirmar el bloqueo.\\[2mm]
\textbf{Pregunta a Operador:} «¿cómo se evita que se imprima una versión distinta de la aprobada?»; conviene averiguar primero qué control existe.
\par}
\end{column}
\end{columns}
\vspace{2.5mm}
{\fontsize{8.5}{10.5}\selectfont\bfseries\color{Accent}
Conclusión: borrador útil para elicitar; ninguna regla adicional se toma como validada por el cliente.\par}
\vspace{1mm}
\end{frame}

\section{CONCLUSIONES}
\begin{frame}[t,shrink=14]{Conclusiones de la prueba en PrintFlow}
\fuentesdiapositiva{Resultados: evidencias/respuesta-ia.txt y evidencias/revision-respuesta.txt}
\vspace{1mm}
\begin{columns}[T,onlytextwidth]
\begin{column}{0.47\textwidth}
\etiqueta{APORTE OBSERVADO}
\vspace{1.5mm}
{\fontsize{10.2}{13}\selectfont
Una ejecución de Gemini propuso \textbf{11 preguntas} y \textbf{9 requisitos candidatos} para explorar aprobaciones y el paso a producción.\par}
\vspace{4mm}
\etiqueta{REVISIÓN NECESARIA}
\vspace{1.5mm}
{\fontsize{10.2}{13}\selectfont
Los \textbf{9 requisitos} incluyen reglas por confirmar y \textbf{5 preguntas} necesitan mejoras de redacción. RC-09 tiene respaldo contextual, aunque no en F01--F05.\par}
\end{column}
\begin{column}{0.47\textwidth}
\etiqueta{CONSULTAS PENDIENTES}
\vspace{1.5mm}
{\fontsize{10.2}{13}\selectfont
Acordar vigencia y unidad de aprobación, roles autorizados, simultaneidad, registro histórico y cambios después de iniciar producción.\par}
\vspace{4mm}
\etiqueta{LÍMITE DEL RESULTADO}
\vspace{1.5mm}
{\fontsize{10.2}{13}\selectfont
Es una prueba documental de una herramienta, no una validación con clientes ni una comparación del rendimiento de distintas IA.\par}
\end{column}
\end{columns}
\vspace{3mm}
{\fontsize{9}{11}\selectfont\bfseries\color{Accent}
Gemini aportó un borrador útil; aprobar requisitos exige contrastarlos con los interesados.\par}
\vspace{1mm}
\end{frame}

\section{USO DE IA}
\begin{frame}[t,shrink=14]{Declaración de uso de IA}
\fuentesdiapositiva{Registros: evidencias/ejecucion.txt y evidencias/revision-respuesta.txt}
\vspace{1mm}
\etiqueta{HERRAMIENTAS Y PROPÓSITO}
\vspace{1.5mm}
{\fontsize{10}{13}\selectfont
\textbf{Gemini, modalidad gratuita} (modelo y versión no disponibles): generación de preguntas y requisitos candidatos con el contexto documental de PrintFlow, en una ejecución del 21/11/2026.\par}
\vspace{3mm}
{\fontsize{10}{13}\selectfont
\textbf{ChatGPT, GPT-6}: apoyo para redactar el cierre, preparar archivos de entrega y realizar una comprobación técnica preliminar.\par}
\vspace{4mm}
\etiqueta{CONTROL Y RESPONSABILIDAD HUMANA}
\vspace{1.5mm}
{\fontsize{10}{13}\selectfont
La respuesta original de Gemini se conservó sin editar y se contrastó con el contexto F01--F05 en la revisión del equipo. Las propuestas permanecen pendientes de validación con los interesados.\par}
\vspace{2.5mm}
{\fontsize{9.5}{12}\selectfont\color{Muted}
El equipo deberá verificar esta declaración, los archivos y la compilación antes de publicar la entrega definitiva.\par}
\vfill
\vspace{1mm}
\end{frame}

\section{BIBLIOGRAFÍA}
\begin{frame}[t]{Referencias}
\fuentesdiapositiva{}
\renewcommand*{\bibfont}{\scriptsize}
\setlength{\bibitemsep}{1.5mm}
\printbibliography[heading=none]
\end{frame}
\end{document}
